\begin{lemma}
\label{editorial-source-lemma-geometrically-integral-any-integral-base-change}
Let $k$ be a field. Let $S$ be a geometrically integral $k$-algebra.
Let $R$ be a $k$-algebra and an integral domain. Then $R \otimes_k S$
is an integral domain.
\end{lemma}

\begin{proof}
Let $F=\operatorname{Frac}(R)$.
The original injection $R\to F$
gives an injection
$$
R\otimes_k S\longrightarrow F\otimes_k S
$$
because tensoring $k$-vector spaces
preserves injections. The receiving
ring is a domain by geometric
integrality of $S/k$. The source
is nonzero: extending $1$ to
bases of the two original
nonzero factors shows
$1\otimes1\neq0$. Hence it
is a domain, as asserted.
The corresponding localization
maps retain the original
denominators $W=R\setminus\{0\}$:
$$
F\otimes_k S\longrightarrow W^{-1}(R\otimes_k S),
\qquad(r/w)\otimes s\longmapsto(r\otimes s)/(w\otimes1),
$$
and the inverse sends
$(\sum_i r_i\otimes s_i)/(w\otimes1)$
to $\sum_i(r_i/w)\otimes s_i$.
The fraction equivalence relations
and tensor balancing make these
well-defined inverses. The
displayed injection is exactly
the original constant-fraction
map under this comparison.

The argument gives precise
ideal consequences for any
original $k$-algebra $R$, now
without a domain hypothesis.
Write $A=R\otimes_k S$ and
let $I\subseteq R$ be an
arbitrary ideal. The maps
$[r\otimes s]\mapsto[r]\otimes s$
and $[r]\otimes s\mapsto[r\otimes s]$
identify
$$
A/IA\simeq(R/I)\otimes_k S.
$$
Choose a $k$-linear functional
$\ell:S\to k$ with $\ell(1)=1$.
The map $\operatorname{id}\otimes\ell$
splits $R/I\to(R/I)\otimes_k S$
as a linear map, proving
$IA\cap R=I$ for every $I$.
If $I$ is prime, $R/I$ is
a domain and the argument
above makes $A/IA$ a domain,
so $IA$ is prime. Conversely,
if $IA$ is prime, its
contraction $I$ is prime.
This proves the exact equivalence
for the original extended ideal,
including the exclusion of
the unit ideal from prime ideals.

There is also the full radical
comparison
$$
\sqrt{IA}=(\sqrt I)A.
$$
To prove it, keep the quotient
map
$$
(R/I)\otimes_k S\longrightarrow(R/\sqrt I)\otimes_k S.
$$
Its kernel is the ideal
generated by the original
$\sqrt I/I$ factor, by the
same explicit quotient maps.
Each element of this kernel
is a finite sum of tensors
$r_j\otimes s_j$ with $r_j$
nilpotent. If $r_j^{d_j}=0$,
its power of degree
$N=1+\sum_j(d_j-1)$ vanishes:
in the full multinomial sum
$$
\sum_{i_1+\cdots+i_m=N}
\frac{N!}{i_1!\cdots i_m!}
\prod_{j=1}^m(r_j\otimes s_j)^{i_j},
$$
some $i_j\geq d_j$ in every
summand. The displayed coefficient
is an integer; no factorial
is inverted in the receiving ring.
The target is reduced by Lemma
\ref{editorial-source-lemma-geometrically-reduced-any-reduced-base-change},
since $S$ is geometrically reduced
and $R/\sqrt I$ is reduced.
Thus this kernel is exactly
the nilradical of $A/IA$.
Taking its inverse image in
$A$ proves the radical formula.

In particular for every prime
$\mathfrak p$ of $R$ the ideal
$\mathfrak pA$ is already prime,
so the continuous section of
Lemma
\ref{editorial-source-lemma-geometrically-irreducible-any-base-change}
has the exact value
$\sigma(\mathfrak p)=\mathfrak pA$.
Its contraction, topology,
base-change compatibility and
minimal-prime correspondence
are those already proved for
the original tensor map;
the stronger integrality
hypothesis removes no terms
from the original ideal.
\end{proof}